\documentclass[11pt]{article}
\usepackage[a4paper,margin=25mm]{geometry}
\usepackage{amsmath,amssymb,amsthm}
\usepackage[hidelinks]{hyperref}
\emergencystretch=2em
\newtheorem{theorem}{Theorem}
\title{The fixed-dimensional lower convexity envelope is zero\\
\large Disproof of Conjecture 00000000947}
\author{gaochengzhecpu}\date{}
\begin{document}\maketitle
\begin{abstract}
Already in dimension two, the lower envelope of the
modulus of convexity over all norms is zero. Every finite
$\ell_q^2$, $1<q<\infty$, has strictly positive modulus
at positive separation, so a minimizing exponent cannot
tend to two as the separation tends to zero. Even if one
restricts the envelope to uniformly convex norms, its
infimum remains zero for separations below two.
Lean verifies actual norm axioms, genuine modulus infima,
strict positivity, the exponent-selection contradiction,
and the uniformly convex approximation.
\end{abstract}

\section*{The genuine modulus and its lower envelope}
Work on the fixed real vector space $E=\mathbb R^2$.
For any norm $N$ on $E$ and $0\leq\varepsilon\leq2$, use
the standard unit-sphere form of Clarkson's modulus
\cite[Section 1]{llorens}:
\[
 \begin{split}
 P_N(\varepsilon)&=\{(x,y):N(x)=N(y)=1,\ N(x-y)\geq\varepsilon\},\\
 \delta_N(\varepsilon)&=
 \inf_{(x,y)\in P_N(\varepsilon)}
       \left(1-N\left(\frac{x+y}{2}\right)\right).
 \end{split}
\]
This is equivalent to the unit-ball version in this
dimension; the cited section recalls the equivalent
conventions. The infimum here is over all admissible
unit pairs, not just symmetric test vectors.
It is nonempty: for any nonzero vector $e$, the pair
$x=e/N(e)$ and $-x$ has separation two.
The triangle inequality makes every displayed defect
nonnegative. Thus $\delta_N(\varepsilon)\geq0$.
Let $L(\varepsilon)=\inf_N\delta_N(\varepsilon)$, over
all genuine norms on this fixed plane.

For the actual sup norm $N_\infty(s,t)=\max(|s|,|t|)$,
take
\[
 u=(1,\varepsilon/2),\qquad v=(1,-\varepsilon/2).
\]
Both have norm one, their distance is $\varepsilon$,
and their midpoint has norm one. Their defect is zero.
Together with nonnegativity this proves
\[
 \delta_{\ell_\infty^2}(\varepsilon)=0,
 \qquad L(\varepsilon)=0\quad(0\leq\varepsilon\leq2).
\]

\newpage
\section*{Positive modulus for every finite exponent greater than one}
For real $1<q<\infty$, use the genuine norm
\[
 N_q(x)=\left(|x_0|^q+|x_1|^q\right)^{1/q}.
\]
We prove positivity of its actual modulus without
assuming an explicit Hanner formula.
For real $a,b$, convexity and the absolute-value triangle
inequality give
\[
 \left|\frac{a+b}{2}\right|^q
 \leq\frac{|a|^q+|b|^q}{2}.
\]
The inequality is strict when $a\ne b$. If $|a|\ne|b|$,
strict convexity of $t\mapsto t^q$ on $[0,\infty)$ gives
strictness after the triangle inequality. If $|a|=|b|$
but $a\ne b$, then $a=-b\ne0$, and strictness is immediate.
Consequently distinct $N_q$-unit vectors satisfy
\[
 N_q\left(\frac{x+y}{2}\right)^q
 <\frac{N_q(x)^q+N_q(y)^q}{2}=1.
\]

Fix $0<\varepsilon\leq2$. The admissible set
$P_{N_q}(\varepsilon)$ is nonempty and closed.
Each coordinate of each unit vector has absolute value
at most one, so this set is bounded in $\mathbb R^4$,
hence compact. The continuous defect attains its
minimum there. Its value is strictly positive at every
admissible pair, since the positive separation forces
$x\ne y$. Thus the attained minimum is positive:
\[
 \delta_{\ell_q^2}(\varepsilon)>0
 \qquad(1<q<\infty,\ 0<\varepsilon\leq2).
\]
Compactness is essential here; pointwise positivity
alone would not justify a positive infimum.

The same conclusions also follow directly for the
unit-ball convention $N_q(x),N_q(y)\leq1$. Its admissible
set is again compact. If either norm is below one, the
triangle inequality makes the midpoint norm strictly
less than one; if both are one, use the strict inequality
above. Thus its attained minimum is positive as well.
The displayed sup-norm pair still gives exact zero.

\begin{theorem}
No choice of minimizing exponents $q(\varepsilon)$
for the fixed-dimensional lower envelope can satisfy
$q(\varepsilon)\to2$ as $\varepsilon\to0^+$.
\end{theorem}
\begin{proof}
Such convergence forces $1<q(\varepsilon)<\infty$ for
all sufficiently small positive $\varepsilon$.
At any such $\varepsilon<2$, its modulus is strictly
positive, while the lower envelope is zero. Therefore
it does not attain that envelope.
\end{proof}

This contradicts the source's asserted small-separation
selection already for dimension two. No assertion about
the separate large-separation limit or endpoint slopes
is needed to refute the stated conjunction.

\newpage
\section*{Restricting to uniformly convex norms does not repair it}
The exact source imposes no uniformly convex norm
restriction. Nevertheless, the same obstruction holds
on that smaller class for $0<\varepsilon<2$.

The positive modulus just proved shows that every finite
$\ell_q^2$, $q>1$, is uniformly convex. For such a $q$ put
$k_q=2^{1/q}$ and take the actual unit vectors
\[
 x_q=\frac{(1,1)}{k_q},\qquad
 y_q=\frac{(1,-1)}{k_q}.
\]
Their separation is $2/k_q$ and their midpoint norm is
$1/k_q$. Hence, whenever $\varepsilon\leq2/k_q$, the
genuine modulus satisfies the \emph{upper bound}
\[
 0<\delta_{\ell_q^2}(\varepsilon)
 \leq1-2^{-1/q}.
\]
We do not assign this trial-pair expression as an exact
modulus formula.

As $q\to\infty$, $2^{-1/q}\to1$. Given any fixed
$0<\varepsilon<2$ and any $\eta>0$, choose a finite
$q\geq2$ large enough that
\[
 2\cdot2^{-1/q}>\varepsilon,
 \qquad 1-2^{-1/q}<\eta.
\]
This produces a genuinely uniformly convex norm with
$0<\delta_{\ell_q^2}(\varepsilon)<\eta$.
Since every norm modulus is nonnegative, even the
infimum over only uniformly convex norms is zero.
Every finite $q>1$ lies strictly above it, so the claimed
minimizing selection tending to two remains impossible.
This approximation assertion deliberately excludes
$\varepsilon=2$; no conclusion about that endpoint is
needed.

\section*{Actual geometric definitions in Lean}
The underlying vector space is \texttt{Fin 2} to
$\mathbb R$. The predicate \texttt{IsNorm} consists
exactly of nonnegativity, positive definiteness,
the triangle inequality and absolute homogeneity.
The definitions \texttt{modulus} and \texttt{envelope}
are actual infima over all admissible unit pairs and
all such norms. Domain nonemptiness is proved for every
genuine norm, so no empty-infimum default produces the
zero envelope.

Finite $\ell_q$ norms are the actual Mathlib
\texttt{PiLp} norms, pulled back by the canonical
coordinate equivalence. Their sum-of-real-powers formula,
continuity, strict midpoint inequality and compact
admissible sets are proved. A general uniformly convex
instance is derived from these facts, rather than
assumed as an unavailable library instance.

\newpage
\section*{Formal conclusions and reproducibility}
The principal theorems are:
\begin{itemize}
\item \texttt{pairs\_nonempty} and
\texttt{envelope\_zero}: all norm moduli are properly
defined on the domain and the exact envelope is zero.
\item \texttt{lp\_modulus\_pos} and
\texttt{no\_finite\_lp\_minimizer}: every finite
$1<q<\infty$ has positive actual modulus and cannot
attain the zero envelope.
\item \texttt{no\_selector\_tending\_two}: direct
negation of an arbitrary minimizing exponent selector
whose limit as $\varepsilon\to0^+$ is two.
\item \texttt{lp\_uniformConvex} and
\texttt{arbitrarily\_small\_uniformly\_convex\_lp}:
actual Mathlib uniform-convexity certificates and
arbitrarily small positive moduli in that class.
\end{itemize}
The selector theorem permits extended exponents as well:
convergence of their real values to two forces them
eventually into the finite range $q>1$. The approximation
theorem gives a real exponent and an actual uniformly
convex \texttt{PiLp} space, together with strict
$0<\delta<\eta$, for every fixed $0<\varepsilon<2$.

Lean 4.19.0 and Mathlib revision\\
\texttt{c44e0c8ee63ca166450922a373c7409c5d26b00b}
are pinned. From the submitted \texttt{lean/} directory:
\begin{verbatim}
lake exe cache get
lake build
lake env lean -DwarningAsError=true Main.lean
\end{verbatim}
The entire project is built in a fresh directory. The
strict check treats warnings as errors, and the eight
principal axiom audits contain only standard foundational
axioms. There are no proof placeholders, custom axioms,
numerical approximations or external computation oracles.
No auxiliary program is required. Compile the document
with \texttt{tectonic main.tex}.

The exact bilingual source, provenance, actual build
logs, artifact hashes and PDF verification records are
included. An agent developed the proof and supplied an
adversarial self-review. The parent separately reviewed
the complete proof, its scope, verification evidence and
every final PDF page. No external independent review or
novelty is claimed for these elementary convexity facts.

\begin{thebibliography}{9}
\bibitem{source} The Last Math Competition, Conjecture
00000000947. Exact bilingual source and provenance included.
\bibitem{llorens} E. Llorens Fuster,
\emph{Moduli and Constants}, April 14, 2006, Section 1.
\href{https://www.uv.es/~llorens/Documento.pdf}{Author-hosted survey}.
The citation supplies standard terminology and equivalent
modulus conventions. All infimum and strict-positivity
arguments used in this disproof are given above.
\end{thebibliography}
\end{document}
